\documentclass[11pt]{article}
\usepackage{progtable}

\begin{document}
\progheader{Quatrième}{4e}

\begin{proglist}

\seq{4C01}{Relatifs — rappels et repérage}{NUM}{Consolider les relatifs acquis en 5e}{
\item 4C01-NUM01 — Comparer et ranger des nombres relatifs
\item 4C01-NUM02 — Additionner et soustraire des relatifs
\item 4C01-NUM03 — Repérer un point dans un plan (x, y)
}

\seq{4C02}{Théorème de Pythagore}{GEO}{Maîtriser Pythagore et les racines carrées}{
\item 4C02-GEO01 — Énoncer et appliquer le théorème de Pythagore
\item 4C02-GEO02 — Calculer l’hypoténuse d’un triangle rectangle
\item 4C02-GEO03 — Calculer une racine carrée
}

\seq{4C03}{Produits et quotients de relatifs}{CALC}{Maîtriser les quatre opérations sur ℤ}{
\item 4C03-CALC01 — Déterminer le signe d’un produit
\item 4C03-CALC02 — Multiplier des nombres relatifs
\item 4C03-CALC03 — Diviser des nombres relatifs
}

\seq{4C04}{Opérations sur les fractions}{NUM}{Opérer sur les fractions avec aisance}{
\item 4C04-NUM04 — Simplifier et comparer des fractions
\item 4C04-NUM05 — Additionner et soustraire des fractions
\item 4C04-NUM06 — Multiplier des fractions
}

\seq{4C05}{Réciproque de Pythagore}{GEO}{Utiliser la réciproque de Pythagore}{
\item 4C05-GEO04 — Rappeler l’égalité de Pythagore
\item 4C05-GEO05 — Appliquer la réciproque pour montrer qu’un triangle est rectangle
\item 4C05-GEO06 — Résoudre des problèmes avec la réciproque
}

\seq{4C06}{Statistiques descriptives}{STAT}{Décrire une série de données}{
\item 4C06-STAT01 — Calculer étendue et fréquence
\item 4C06-STAT02 — Calculer la moyenne d’une série
\item 4C06-STAT03 — Déterminer la médiane
}

\seq{4C07}{Puissances}{NUM}{Maîtriser le calcul avec les puissances}{
\item 4C07-NUM07 — Calculer des produits et quotients de puissances
\item 4C07-NUM08 — Calculer des puissances de puissances
\item 4C07-NUM09 — Utiliser les puissances de 10 et l’écriture décimale
}

\seq{4C08}{Proportionnalité}{PROP}{Approfondir le raisonnement proportionnel}{
\item 4C08-PROP01 — Reconnaître une situation proportionnelle
\item 4C08-PROP02 — Utiliser coefficient et produit en croix
\item 4C08-PROP03 — Résoudre des problèmes de proportionnalité
}

\seq{4C09}{Périmètres et aires}{MEAS}{Généraliser le calcul d’aires et périmètres}{
\item 4C09-MEAS01 — Calculer le périmètre de figures composées
\item 4C09-MEAS02 — Calculer l’aire de figures composées
\item 4C09-MEAS03 — Choisir la bonne formule selon la figure
}

\seq{4C10}{Sommes de fractions (dénom. différents)}{NUM}{Consolider les sommes de fractions}{
\item 4C10-NUM10 — Décomposer et réduire au même dénominateur
\item 4C10-NUM11 — Additionner des fractions à dénominateurs multiples
\item 4C10-NUM12 — Additionner des fractions relatives
}

\seq{4C11}{Calcul littéral}{ALG}{Introduire l’algèbre formelle}{
\item 4C11-ALG01 — Substituer une valeur dans une expression littérale
\item 4C11-ALG02 — Utiliser des programmes de calcul
\item 4C11-ALG03 — Réduire et ordonner une expression algébrique
}

\seq{4C12}{Probabilités}{STAT}{Approfondir le calcul des probabilités}{
\item 4C12-STAT04 — Déterminer les issues possibles d’une expérience aléatoire
\item 4C12-STAT05 — Calculer une probabilité par dénombrement
\item 4C12-STAT06 — Utiliser la notion d’événement contraire
}

\end{proglist}
\end{document}
